% Generated draft: requires mathematical and source review.
\documentclass[11pt]{article}
\usepackage[margin=25mm]{geometry}
\usepackage{fontspec}
\setmainfont{Latin Modern Roman}
\usepackage{amsmath,amssymb,mathtools,unicode-math}
\setmathfont{Latin Modern Math}
\usepackage{enumitem,needspace,xurl,hyperref}
\setlength{\parindent}{0pt}
\setlength{\parskip}{5pt}
\setlength{\emergencystretch}{3em}
\newcommand{\N}{\mathbb{N}}
\newcommand{\Z}{\mathbb{Z}}
\newcommand{\Q}{\mathbb{Q}}
\newcommand{\R}{\mathbb{R}}
\newcommand{\C}{\mathbb{C}}
\newcommand{\F}{\mathbb{F}}
\newcommand{\A}{\mathbb{A}}
\newcommand{\PP}{\mathbb{P}}
\newcommand{\E}{\mathbb{E}}
\newcommand{\eps}{\varepsilon}
\newcommand{\dd}{\,\mathrm{d}}
\begin{document}
\section*{UnsolvedMath ID: 20002451}
\textbf{Source problem code:} catalogue entry
\subsection{Definitions and mathematical statement}
The thin-film question, identified by the supplied verification as P. Howard's, concerns
\[
 u_t+f(u)_x=-\eps(u^3u_{xxx})_x,\qquad f(u)=u^2-u^3,\quad \eps>0,
\]
and the stability or instability of a separating combination of a left-moving Lax shock and a right-moving undercompressive shock. The extract subsequently includes C. Jones's energized-NLS questions and T. Kapitula's Gross--Pitaevskii vortex discussion; these are adjacent entries, not extra terms or hypotheses of this equation. The original question does not specify the selected profiles, perturbation space, or meaning of stability.

The aid's algebraic partial result uses ordered states $0<b<a<m<1$, with a first front $a\to m$ and second front $m\to b$. Rankine--Hugoniot speeds are
\[
 s_1=\frac{f(m)-f(a)}{m-a},\qquad s_2=\frac{f(b)-f(m)}{b-m}.
\]
The first is Lax when $f'(a)>s_1>f'(m)$. The strong-undercompressive convention here is $s_2>\max\{f'(m),f'(b)\}$. The aid states that these inequalities together with $s_1<0<s_2$ hold exactly in the chamber
\[
 b<\tfrac16,\qquad \tfrac23<m<1-2b,\qquad
 T(m)<a<m,\quad
 T(m)=\frac{1-m+\sqrt{(1-m)(1+3m)}}2,
\]
within the ordered-state assumptions. It also gives the exact identity
\[
 s_2-s_1=(a-b)(a+b+m-1).
\]
Thus under the strict ordering the algebraic equal-speed threshold is $a=1-m-b$; the opposite-moving chamber itself has strictly positive separation speed. These algebraic conditions do not prove existence or selection of fourth-order traveling profiles.

Assume selected profiles $U_1,U_2$ exist, have the displayed endstates, and have exponential tails with sufficient derivative bounds for the equation. Their superposition is
\[
 U^{\mathrm{comp}}(t,x)=U_1(x-s_1t-h_1)+U_2(x-s_2t-h_2)-m,
 \qquad d(t)=h_2-h_1+(s_2-s_1)t.
\]
For initial separation $d(0)=d_0$ in the tail-interaction regime, the aid reports a residual estimate, in each $W^{-1,p}$ space covered by those profile bounds,
\[
 \|\partial_tU^{\mathrm{comp}}+\partial_xf(U^{\mathrm{comp}})
 +\eps\partial_x((U^{\mathrm{comp}})^3\partial_x^3U^{\mathrm{comp}})\|_{W^{-1,p}}
 \le C e^{-\kappa d(t)},\qquad
 \int_0^\infty\|\text{residual}\|_{W^{-1,p}}\,\dd t
 \le\frac{C e^{-\kappa d_0}}{\kappa(s_2-s_1)}.
\]
The exponent range and exact tail hypotheses are not specified in the summary. For integrable perturbations, phase shifts contribute mass
\[
 \Delta M=(a-m)\delta h_1+(m-b)\delta h_2;
\]
a remainder contributes its own integral. Hence conserved mass imposes one linear constraint on two phases, not two independent phase conditions. These facts reduce the stability question to a two-phase, nonautonomous linear and modulation problem; the aid does not solve that problem.
\subsection{Short English statement}
The target is the stability of two thin-film fronts moving apart. The supplied partial result identifies the exact algebraic state region, their separation speed, and an integrable interaction error conditional on existing exponentially localized profiles. It leaves the nonlinear stability theorem open.
\subsection{Sources}
\begin{itemize}
\item Local UnsolvedMath registry, permanent ID 20002451, source code catalogue entry; created and updated 14 August 2026. Canonical \texttt{aim-pdes-notes.json}, zero-based \texttt{notes[62]}, tag ``open problem,'' item 15.
\item AIM workshop report ``Open Problems in Multidimensional Stability of Waves and Patterns,'' P. Howard problem as identified by the supplied verification: \url{https://aimath.org/WWN/multidimwaves/multidimwaves.pdf}. The ``Combination Structure'' figure is referenced but not supplied. The following C. Jones and T. Kapitula material is record-boundary contamination. Original remarks and literature field are empty.
\item Research attempt 1: \texttt{partial\_result}, problem status at run open, rigor label \texttt{proved}, novelty confidence low. No complete authoritative publication reference or full profile/stability theorem is supplied for the lemma package.
\item Ulam AI / UnsolvedMath Contributors, CC BY 4.0. Machine-generated research aid, not peer reviewed; independently verify before citation.
\end{itemize}
\end{document}
